% Shared visual style for all lecture-note projects.

% Document color palette
\definecolor{ChapterBlue}{HTML}{12355B}
\definecolor{SectionBlue}{HTML}{1F5F99}
\definecolor{SubsectionBlue}{HTML}{4F7FA8}
\definecolor{AccentBlue}{HTML}{D8E9F7}
\definecolor{ExampleGreen}{HTML}{2F7D4A}
\definecolor{ExampleGreenBack}{HTML}{EEF8F1}
\definecolor{TheoremOrange}{HTML}{B85C00}
\definecolor{TheoremOrangeBack}{HTML}{FFF4E6}

% Chapter, section, and subsection heading styles
\titleformat{\chapter}[display]
    {\normalfont\huge\bfseries\color{ChapterBlue}}
    {\chaptertitlename\ \thechapter}
    {0.6em}
    {}

\titlespacing*{\chapter}{0pt}{-0.5em}{2.0em}

\titleformat{\section}
    {\normalfont\Large\bfseries\color{SectionBlue}}
    {\thesection}
    {0.8em}
    {}

\titlespacing*{\section}{0pt}{1.8em}{0.8em}

\titleformat{\subsection}
    {\normalfont\large\bfseries\color{SubsectionBlue}}
    {\thesubsection}
    {0.8em}
    {}

\titlespacing*{\subsection}{0pt}{1.2em}{0.5em}

% Theorem-like environments
\newtheorem{theorem}{Theorem}[chapter]
\newtheorem{proposition}{Proposition}[chapter]
\newtheorem{lemma}{Lemma}[chapter]
\newtheorem{corollary}{Corollary}[chapter]

\theoremstyle{definition}
\newtheorem{definition}{Definition}[chapter]
\newtheorem{example}{Example}[chapter]

\theoremstyle{remark}
\newtheorem{remark}{Remark}[chapter]

% Colored boxes
\newtcolorbox{definitionbox}{
    breakable,
    colback=AccentBlue!28,
    colframe=SectionBlue,
    title=Definition,
    fonttitle=\bfseries
}

\newtcolorbox{theorembox}{
    breakable,
    colback=TheoremOrangeBack,
    colframe=TheoremOrange,
    title=Theorem,
    fonttitle=\bfseries
}

\newtcolorbox{examplebox}{
    breakable,
    colback=ExampleGreenBack,
    colframe=ExampleGreen,
    title=Example,
    fonttitle=\bfseries
}

\newtcolorbox{remarkbox}{
    breakable,
    colback=gray!4,
    colframe=SubsectionBlue!75,
    title=Remark,
    fonttitle=\bfseries
}

\newtcolorbox{summarybox}{
    breakable,
    colback=gray!5,
    colframe=ChapterBlue,
    title=Summary,
    fonttitle=\bfseries
}


% Project-specific Polish box titles and practical-lab environments.
\renewtcolorbox{definitionbox}{
    breakable,
    colback=AccentBlue!28,
    colframe=SectionBlue,
    title=Definicja,
    fonttitle=\bfseries
}

\renewtcolorbox{theorembox}{
    breakable,
    colback=TheoremOrangeBack,
    colframe=TheoremOrange,
    title=Twierdzenie,
    fonttitle=\bfseries
}

\renewtcolorbox{examplebox}{
    breakable,
    colback=ExampleGreenBack,
    colframe=ExampleGreen,
    title=Przykład,
    fonttitle=\bfseries
}

\renewtcolorbox{remarkbox}{
    breakable,
    colback=gray!4,
    colframe=SubsectionBlue!75,
    title=Uwaga,
    fonttitle=\bfseries
}

\renewtcolorbox{summarybox}{
    breakable,
    colback=gray!5,
    colframe=ChapterBlue,
    title=Podsumowanie,
    fonttitle=\bfseries
}

\definecolor{WarningRed}{HTML}{A33A3A}
\definecolor{WarningRedBack}{HTML}{FFF2F2}

\newtcolorbox{warningbox}{
    breakable,
    colback=WarningRedBack,
    colframe=WarningRed,
    title=Typowy błąd,
    fonttitle=\bfseries
}

\newtcolorbox{exercisebox}{
    breakable,
    colback=AccentBlue!18,
    colframe=SectionBlue!85,
    title=Ćwiczenie laboratoryjne,
    fonttitle=\bfseries
}

% Project-specific styling for code examples.
% The notes use many standard verbatim blocks for Mathematica code.
% Redefine verbatim through fvextra so code is not printed on a stark white background.
\usepackage{fvextra}

\definecolor{CodeGrayBack}{HTML}{F7F7F7}
\definecolor{CodeGrayFrame}{HTML}{E2E2E2}

\DefineVerbatimEnvironment{verbatim}{Verbatim}{
    breaklines=true,
    breakanywhere=true,
    fontsize=\small,
    frame=single,
    framerule=0.25pt,
    rulecolor=\color{CodeGrayFrame},
    framesep=3mm,
    xleftmargin=0pt,
    xrightmargin=0pt,
    bgcolor=CodeGrayBack
}
